\section{总结与延伸}

本节把 Lovart 特辑压缩成几个结论。第一，Lovart 是垂直 Agent，不是通用 Agent；它的场景是设计师工作流。第二，AI 应用创业的生存线很残酷，补贴战、下架、现金见底和融资困难都会改变产品判断。第三，创始人的“爽”是 PMF 情绪信号，但必须转成留存、付费和工作流。第四，应用公司不能只依赖模型红利，要沉淀垂直上下文和产品心智。

\begin{importantbox}{把 Lovart 特辑放进张小珺 AI/互联网队列}
EP110 讲 Agent 技术报告，EP112 讲 AI 产品像挖矿，EP115 讲 Agent 下半场理论；Lovart 特辑则是一个应用创业现场切片：当模型能力窗口打开，一个垂直 Agent 如何从低谷走到全球关注。
\end{importantbox}

\begin{knowledgebox}{与相关节目的关系}
\begin{tabularx}{\textwidth}{p{0.18\textwidth}p{0.34\textwidth}X}
\toprule
节目 & 主题 & 与 Lovart 的连接 \\
\midrule
EP112 & AI 产品像挖矿，窗口期变短 & Lovart 是窗口期里的垂直应用案例。 \\
EP110 & Agent 技术报告和上下文工程 & 解释垂直 Agent 背后的工具/上下文能力。 \\
EP128 & Manus 出售前访谈 & 对照通用 Agent 与垂直 Agent 的不同路径。 \\
EP113 & Kimi K2 和 Agentic LLM & 解释模型能力窗口如何打开应用机会。 \\
\bottomrule
\end{tabularx}
\end{knowledgebox}

\subsection{关键 takeaways}

\begin{enumerate}
\item 垂直 Agent 的价值在专业工作流，不在“什么都能做”。
\item 设计 Agent 的核心是 brief、素材、审美、版本和交付。
\item AI 应用的 Magic Moment 必须快速转成留存和付费。
\item 创业低谷是产品判断的一部分，不只是故事背景。
\item 通用 Agent 和垂直 Agent 可能并存，关键在入口与深度的分工。
\end{enumerate}

\subsection{开放问题}

前面已经总结了 Lovart 的产品逻辑和创业路径，本节保留几个仍然需要继续观察的问题。这些问题决定 Lovart 是成为短期爆款，还是成为长期设计工作台。

\begin{enumerate}
\item Lovart 的设计工作流壁垒能否抵御通用 Agent 的下探？
\item 设计师是否愿意把关键审美判断交给 Agent？
\item AI 设计工具的商业模式更像订阅、按项目收费，还是团队工作台？
\item 垂直 Agent 如何处理版权、素材来源和品牌一致性？
\end{enumerate}

\subsection{拓展阅读}

\begin{itemize}
\item 对 Agent 技术底层感兴趣，可对照 EP110 技术报告串讲。
\item 对 AI 产品窗口感兴趣，可对照 EP112 大模型季报。
\item 对 Agent 创业组织感兴趣，可对照 EP128 Manus 和 EP113 Kimi K2 访谈。
\end{itemize}

\end{document}
